% Part: incompleteness
% Chapter: theories-computability
% Section: consis-with-q

\documentclass[../../../include/open-logic-section]{subfiles}

\begin{document}

\olfileid{inc}{tcp}{con}

\olsection{Teori kang konsisten karo $\Th{Q}$ ora bisa diputusake}

Teorema sabanjure nyatakake yen ora mung $\Th{Q}$ kang ora bisa
diputusake, nanging sejatine saben teori kang ora pasulayan karo
$\Th{Q}$ uga ora bisa diputusake.

\begin{thm}
Upamana $\Th{T}$ teori sembarang ing basa aritmetika kang konsisten
karo $\Th{Q}$ (yaiku, $\Th{T} \cup \Th{Q}$ konsisten). Mula
$\Th{T}$ ora bisa diputusake.
\end{thm}

\begin{proof}
Elinga yen $\Th{Q}$ nduweni himpunan aksioma winates, $!Q_1$, \dots,
$!Q_8$. Kita malah bisa ngganti kabeh iki nganggo siji aksioma,
$!E = !Q_1 \land \dots \land !Q_8$.

Upamana $\Th{T}$ teori kang bisa diputusake lan konsisten karo~$\Th{Q}$.
Tetepna
\[
C = \Setabs{!A}{\Th{T} \Proves !E \lif !A}.
\]
Kita tuduhake yen $C$ bakal dadi pamisah komputabel kanggo $\Th{Q}$ lan
$\Th{\bar Q}$, kang dadi kontradiksi. Sepisan, yen $!A$ kalebu ing
$\Th{Q}$, mula $!A$ bisa dibuktekake saka aksioma $\Th{Q}$; miturut
teorema deduksi, ana !!a{derivation} kanggo $!E \lif !A$ ing logika orde siji.
Mula $!A$ kalebu ing~$C$.

Kosok baline, yen $!A$ kalebu ing $\Th{\bar Q}$, mula ana bukti
kanggo $!E \lif \lnot !A$ ing logika orde siji. Yen $\Th{T}$ uga
mbuktekake $!E \lif !A$, mula $\Th{T}$ mbuktekake $\lnot !E$, lan ing
kahanan iku $\Th{T} \cup \Th{Q}$ ora konsisten. Nanging kita nganggep
yen $\Th{T} \cup \Th{Q}$ konsisten, mula $\Th{T}$ ora mbuktekake
$!E \lif !A$, lan $!A$ ora kalebu ing~$C$.

Kita wis nuduhake: yen $!A$ kalebu ing $\Th{Q}$, mula uga kalebu ing
$C$; lan yen $!A$ kalebu ing $\Th{\bar Q}$, mula kalebu ing
$\Complement{C}$. Dadi $C$ iku pamisah komputabel, kang dadi
kontradiksi kang digoleki.
\end{proof}

Teorema iki kuwat banget. Salah siji akibate yaiku:

\begin{cor}
  Logika orde siji kanggo basa aritmetika (yaiku, himpunan
  $\Setabs{!A}{\text{$!A$ bisa dibuktekake ing logika orde siji}}$)
  ora bisa diputusake.
\end{cor}

\begin{proof}
Logika orde siji iku himpunan akibat saka $\emptyset$,
kang konsisten karo~$\Th{Q}$.
\end{proof}

\end{document}
